\documentclass{book}
\usepackage[utf8]{inputenc}
\usepackage[T1]{fontenc}
\usepackage{geometry}
\usepackage{graphicx}
\usepackage{float}
\usepackage{hyperref}
\usepackage[english,indonesian]{babel}

\title{Optimasi Pemrosesan Template LaTeX Menggunakan Pendekatan Berbasis AI}
\author{Academic Support Agent}
\date{\today}

\begin{document}
\maketitle

\begin{abstract}

\end{abstract}

\chapter{PENDAHULUAN}


\chapter{TINJAUAN PUSTAKA}


% Original content (commented out):
% \documentclass[12pt, oneside]{book}
\pdfsuppresswarningpagegroup=1

% \listfiles

\usepackage[utf8]{inputenc} %fuentes
\usepackage{lmodern} %fuentes
\usepackage[T1]{fontenc} %fuentes
\usepackage[spanish,english]{babel} %idoma, cargué los dos porque constantemente estoy cambiando
\usepackage{textgreek} %escribir texto en griego sin estar entre $ $
\usepackage{amsmath}
\usepackage{amsthm, amssymb, amsfonts}
\usepackage{graphicx, color}
\usepackage{caption, subcaption} %figuras y subfiguras
\usepackage{multicol} 
\usepackage{makeidx} %Index alfabético
\makeindex
\usepackage{array, tabularx, multirow} %herramientas para tablas
\usepackage[width=155mm,top=25mm,bottom=25mm, headheight=15pt,
]{geometry}
% \usepackage[displaymath,mathlines, pagewise]{lineno}
% \linenumbers
\usepackage[pagebackref=true]{hyperref} %links
\usepackage{enumitem} %herramientas para listas
\hypersetup{
	hidelinks,
  colorlinks  = true,    % Colours links instead of ugly boxes
  urlcolor    = red,    % Colour for external hyperlinks
  linkcolor   = blue,    % Colour of internal links
  citecolor   = blue,      % Colour of citations
  pdfauthor   = {Pito Pérez},%
  pdfsubject  = {Ph. D. Thesis},%
  pdftitle    = {Un ensayo sobre su inútil vida},  %
,}


\usepackage{tikz-cd} %Diagramas conmutativos



\usepackage{lipsum}  %blindText not really necesary

\usepackage[capitalise,noabbrev, nameinlink]{cleveref} %CleverRef


%CleverRef config
\Crefname{part}{Part}{Parts}
% \crefname{prope}{property}{properties}
% \Crefname{part}{Part}{Parts}
\Crefname{step}{Step}{Steps}
\Crefname{prop}{Proposition}{Propositions}
\Crefname{prob}{Problem}{Problems}



\usepackage{url}%urls en bibliografia

\usepackage{fancyhdr} %páginas bonitas

\usepackage{titling} %titlepage
\usepackage{datetime}%titlepage

% \usepackage{breqn}


\graphicspath{ {img/} } %Aquí pondría mis figuras¡Si tuviera unas!


% ************Page Layout **************************%

\pagestyle{fancy}
\fancyhf{} %Remove content 
%\fancyhead[LE,RO]{\bfseries\thepage} %two sides
%\fancyhead[RE]{\nouppercase{\bfseries\leftmark}} %two sides
\fancyhead[LO]{\nouppercase{\bfseries\rightmark}}
\renewcommand{\headrulewidth}{0.5pt}
\renewcommand{\footrulewidth}{0pt}
\addtolength{\headheight}{2pt} 



\makeatletter
  \def\cleardoublepage{\clearpage\if@twoside \ifodd\c@page\else
  \hbox{}
  \vspace*{\fill}
  \begin{center}
  %% "page intentionaly left empty"
  \end{center}
  \vspace{\fill}
  \thispagestyle{empty}
  \newpage
  \if@twocolumn\hbox{}\newpage\fi\fi\fi}
\makeatother

%*********************** Theorems **********************************%
\numberwithin{equation}{section}


% \theoremstyle{plain}
% \newtheorem{thm}{Theorem}[section]
% \newtheorem{lem}[thm]{Lemma}
% \newtheorem{prop}[thm]{Proposition}
% \newtheorem{coro}[thm]{Corollary}
% 
% \theoremstyle{definition}
% \newtheorem{defn}[thm]{Definition}
% \newtheorem{exam}[thm]{Example}
% \newtheorem{prob}[thm]{Problem}
% 
% \theoremstyle{remark}
% \newtheorem{note}[thm]{Note}
% \newtheorem{rem}[thm]{Remark}



\theoremstyle{plain}
\newtheorem{thm}[equation]{Theorem}
\newtheorem{lem}[equation]{Lemma}
\newtheorem{prop}[equation]{Proposition}
\newtheorem{coro}[equation]{Corollary}

\theoremstyle{definition}
\newtheorem{defn}[equation]{Definition}
\newtheorem{exam}[equation]{Example}
\newtheorem{prob}[equation]{Problem}


\theoremstyle{remark}
\newtheorem{note}[equation]{Note}
\newtheorem{rem}[equation]{Remark}



%*********************** User commands ****************************%
%letters
\newcommand{\cP}{\mathcal{P}}

%operators
\DeclareMathOperator{\aut}{Aut}


%shortcuts
\newcommand{\scatter}{{scattering}}


%lists
\renewcommand{\labelitemi}{\tiny{$\blacksquare$}}
\renewcommand{\labelitemii}{$\cdot$}


%others
\newcommand{\sm}{\setminus} %short setminus 
\renewcommand{\leq}{\leqslant}
\renewcommand{\geq}{\geqslant}
\renewcommand{\subset}{\subseteq}
\renewcommand{\qedsymbol}{$\blacksquare$} %qed
\renewcommand{\epsilon}{\varepsilon}
\renewcommand{\bar}{\overline}
% \renewcommand{\cong}{\simeq}

\usepackage[shadow, colorinlistoftodos,disable]{todonotes}
%\setuptodonotes{fancyline, color=blue!30}


%*********************** ************ ****************************%



\begin{document}


\author{Pito Pérez}
\title{Un largo ensayo sobre su inútil vida}
%\date{\monthname, \the\year}
\date{Agosto, 2019}


\frontmatter
\input{chapters/portada}
\cleardoublepage
\include{chapters/dedi}
\cleardoublepage
\include{chapters/dinero}
\cleardoublepage
\renewcommand{\chaptermark}[1]{\markboth{#1}{}}
\renewcommand{\sectionmark}[1]{\markright{#1}}
\selectlanguage{english}
\tableofcontents
% \cleardoublepage
% \listoftodos
\addcontentsline{toc}{chapter}{Resumen/Abstract}
\include{chapters/abstract}
\cleardoublepage
\selectlanguage{spanish}
\include{chapters/acknowlmts}
\cleardoublepage
\selectlanguage{spanish}
\include{chapters/intro_es}
\cleardoublepage
\selectlanguage{english}
\include{chapters/intro}
% 
% \renewcommand{\chaptermark}[1]{\markboth{\thechapter. #1}{}}
% \renewcommand{\sectionmark}[1]{\markright{\thesection\ #1}}
% %\tableofcontents
\mainmatter
\input{chapters/chap1}
\input{chapters/chap2}
\input{chapters/chap3}
% \include{conc}

% \renewcommand{\chaptermark}[1]{\markboth{#1}{}}
% \renewcommand{\sectionmark}[1]{\markright{#1}}
\backmatter

%\input{app_not}
%\clearpage
\cleardoublepage
\phantomsection
\addcontentsline{toc}{chapter}{Index}
\printindex
% 

\cleardoublepage
\phantomsection
\addcontentsline{toc}{chapter}{Bibliography}

\bibliographystyle{plainurl}
\bibliography{chapters/phDThesis}
% \nocite{*}

\end{document}


\end{document}